\documentclass{travelbook}

\usepackage{mwe}

\title{A Week Elsewhere}
\author{Francois}
\date{Summer 2025}

\begin{document}

\makebookcover{example-image}{A Week Elsewhere}{A small book of sun, streets, and sea}

\thispagestyle{empty}
\null
\clearpage

\pagenumbering{arabic}
\setcounter{page}{1}
\tableofcontents*
\clearpage

\section{Arrival}

\chapterspread{First Light}{Markets, ferries, and the first long walk into town}{example-image-duck}

The first morning set the pace for the rest of the trip. We left without a
strict plan, following the shade from one narrow street to the next and stopping
whenever a doorway, a sign, or a view across the water asked for a photograph.

By noon the city felt less like a checklist and more like a route we were
learning by memory. The best pictures came from those small pauses: waiting for
the ferry, crossing a square twice, or turning around just after passing a
market stall.

\begin{journey}[Route]
  \journeystep{Station}{Old Town}{
    We started with the short walk from the station into the older streets,
    letting the first signs and storefronts decide the direction.
  }
  \journeystep{Old Town}{Harbor}{
    The route opened suddenly at the water, where the light was stronger and
    the pace slowed down around the ferry terminal.
  }
  \journeystep{Fish market}{}{
    A short stop between the larger movements of the day, useful as a marker
    without turning it into a full route segment.
  }
  \journeystep{Harbor}{Hilltop}{
    A late climb gave the day its last view: rooftops below, boats moving out,
    and the map finally making sense from above.
  }
\end{journey}

\begin{travelnote}[Arrival]
The best travel days leave a few details out of focus: the first coffee, the
shape of the streets, and the moment the map starts to make sense.
\end{travelnote}

\fullpagephoto[Morning on the waterfront][height=\paperheight]{example-image-a}

\twophotos
  {example-image-b}
  {A quiet street before lunch.}
  {example-image-c}
  {Late afternoon color near the old harbor.}

\section{Harbor Days}

The harbor became the place we returned to whenever the day needed a reset.
There was always movement there: boats arriving, tables being set, children
leaning over the edge to watch the water. It made a natural divider between the
busier streets inland and the slower rhythm near the sea.

These pages are a mix of wide views and small details. In a real book, this is a
good place for a short sequence that connects several days without explaining
every stop.

\begin{travelnote}[A small pause]
Use short notes between image sequences when the book needs rhythm. The class
keeps them visually quiet so the photographs remain the main event.
\end{travelnote}

\clearpage
\thispagestyle{travelbook}
\begin{tikzpicture}[remember picture,overlay]
  \node[anchor=north west,inner sep=0pt]
    at ([xshift=12mm,yshift=-15mm]current page.north west)
    {\includegraphics[width=124mm]{example-image-16x9}};

  \node[anchor=north west,inner sep=0pt]
    at ([xshift=12mm,yshift=-94mm]current page.north west)
    {\includegraphics[height=87mm]{example-image-9x16}};

  \node[anchor=north west,text width=54mm,align=left]
    at ([xshift=78mm,yshift=-96mm]current page.north west)
    {{\sffamily\small\color{tbmuted}Mixed orientation page\par}
     \vspace{3mm}
     The landscape image gives the scene room to breathe, while the portrait
     image keeps the memory close: the walk back along the water, tables being
     cleared, and the last light on the harbor buildings.};
\end{tikzpicture}
\null

\clearpage
\thispagestyle{travelbook}
\begin{tikzpicture}[remember picture,overlay]
  \node[anchor=north west,inner sep=0pt]
    at ([xshift=12mm,yshift=-15mm]current page.north west)
    {\includegraphics[width=124mm]{example-image-16x9}};

  \node[anchor=north west,text width=54mm,align=left]
    at ([xshift=12mm,yshift=-96mm]current page.north west)
    {{\sffamily\small\bfseries\color{tbsecondary}Mirrored mixed page\par}
     \vspace{3mm}
     This version gives the text the first position in the lower row, then lets
     the portrait image close the page. It is useful when the note should read
     before the smaller visual detail.};

  \node[anchor=north west,inner sep=0pt]
    at ([xshift=78mm,yshift=-94mm]current page.north west)
    {\includegraphics[height=87mm]{example-image-9x16}};
\end{tikzpicture}
\null

\clearpage

\threephotos
  {example-image}
  {A wider editorial image with room for the eye to settle.}
  {example-image-a}
  {Texture and detail.}
  {example-image-b}
  {A second fragment from the same walk.}

\fullpagephoto[The last swim before packing][width=\paperwidth]{example-image-c}

\makebookbackcover{example-image}{Printed in A5 portrait. Replace these placeholder images with your vacation photographs.}

\end{document}
